\section{Conclusion}\label{sec:literature-review-conclusion}
This review highlights that while systems of twinned systems research has advanced considerably in concept and scope, its practical realization remains uneven.
Architectural maturity, standardization, and empirical validation lag behind the rapid conceptual integration of DT and SoS principles. Systems of twinned systems are often designed as isolated prototypes rather than cohesive, dependable infrastructures capable of operating under realistic conditions. Across the literature, the need for greater architectural alignment, explicit dependability mechanisms, and measurable system assurance manifest as a consistent theme. Future progress depends on ensuring reliable and adaptable systems of twinned systems, starting from design rather than assuming these qualities as natural outcomes of integration. These findings establish a foundation for the subsequent chapter, which draws on these gaps to define the requirements for dependable systems of twinned systems architectures. 